\section{The LGV lemma: weighted and generalized versions}

\subsection{The weighted version}

\begin{theorem}[LGV lemma, lattice weight version]
\label{thm.lgv.kpaths.wt}
\lean{LGV.lgv_weighted_lattice}
\leantarget
\leanok
Let $K$ be a commutative ring.

For each arc $a$ of the digraph $\mathbb{Z}^{2}$, let $w(a)$ be
an element of $K$. We call $w(a)$ the \emph{weight} of $a$.

For each path $p$ of $\mathbb{Z}^{2}$, define the \emph{weight} $w(p)$
of $p$ by
\[
  w(p) := \prod_{a \text{ is an arc of } p} w(a).
\]

For each path tuple $\mathbf{p} = (p_1, p_2, \ldots, p_k)$, define the
\emph{weight} $w(\mathbf{p})$ of $\mathbf{p}$ by
\[
  w(\mathbf{p}) := w(p_1)\, w(p_2) \cdots w(p_k).
\]

Let $k \in \mathbb{N}$.  Let
$\mathbf{A} = (A_1, A_2, \ldots, A_k)$ and
$\mathbf{B} = (B_1, B_2, \ldots, B_k)$ be two $k$-vertices.  Then,
\[
  \det\!\left(\left(\sum_{p : A_i \to B_j} w(p)\right)_{1 \le i \le k,\; 1 \le j \le k}\right)
  = \sum_{\sigma \in S_k} (-1)^{\sigma}
    \sum_{\substack{\mathbf{p} \text{ is a nipat} \\ \text{from } \mathbf{A} \text{ to } \sigma(\mathbf{B})}}
    w(\mathbf{p}).
\]
Here, ``$p : A_i \to B_j$'' means ``$p$ is a path from $A_i$ to $B_j$''.
\end{theorem}

\begin{proof}
This is the special case of Theorem~\ref{thm.lgv.kpaths.wt-dg} for the
integer lattice $\mathbb{Z}^2$, which is a path-finite acyclic digraph.
\end{proof}

\subsection{Generalization to acyclic digraphs}

\begin{theorem}[LGV lemma, digraph weight version]
\label{thm.lgv.kpaths.wt-dg}
\lean{LGV.lgv_weighted_digraph}
\leantarget
\leanok
Let $K$ be a commutative ring.

Let $D$ be a path-finite (but possibly infinite) acyclic digraph.

For each arc $a$ of $D$, let $w(a) \in K$. For each path $p$ of $D$,
define $w(p) := \prod_{a \text{ is an arc of } p} w(a)$.
For each path tuple $\mathbf{p} = (p_1, \ldots, p_k)$, define
$w(\mathbf{p}) := w(p_1) \cdots w(p_k)$.

Let $k \in \mathbb{N}$.  Let
$\mathbf{A} = (A_1, \ldots, A_k)$ and
$\mathbf{B} = (B_1, \ldots, B_k)$ be two $k$-tuples of vertices of $D$.
Then,
\[
  \det\!\left(\left(\sum_{p : A_i \to B_j} w(p)\right)_{1 \le i \le k,\; 1 \le j \le k}\right)
  = \sum_{\sigma \in S_k} (-1)^{\sigma}
    \sum_{\substack{\mathbf{p} \text{ is a nipat} \\ \text{from } \mathbf{A} \text{ to } \sigma(\mathbf{B})}}
    w(\mathbf{p}).
\]
\end{theorem}

\begin{proof}
The same argument as for
Proposition~\ref{prop.lgv.kpaths.count} can be used here; just replace
$\operatorname{sign}(\sigma, \mathbf{p}) := (-1)^{\sigma}$ by
$\operatorname{sign}(\sigma, \mathbf{p}) = (-1)^{\sigma} \cdot w(\mathbf{p})$.
The only new observation required is that when we exchange the tails of
two paths in our path tuple, the weight of the path tuple does not change.
This is clear: the weight of a path tuple is the product of the weights of
all arcs in all paths of the tuple; when we exchange the tails of two paths,
some arcs get moved from one path to the other, but the total product stays
unchanged.

More precisely, the proof expands the determinant via the Leibniz formula.
Each summand $(-1)^\sigma \prod_i M_{i,\sigma(i)}$ expands (by
distributivity) into a sum over all path tuples from $\mathbf{A}$ to
$\sigma(\mathbf{B})$, weighted by $(-1)^\sigma \, w(\mathbf{p})$.  Split
the sum into non-intersecting and intersecting path tuples.

For the intersecting path tuples, a sign-reversing involution
$\varphi$ (Lemma~\ref{lem.lgv.signReversing.involutive}) is defined as
follows: given an intersecting path tuple paired with a permutation
$(\sigma, \mathbf{p})$, find the canonical first intersection point,
exchange the tails of the two relevant paths at that point, and compose
$\sigma$ with the corresponding transposition.  This involution
reverses the sign factor $(-1)^\sigma$ (Lemma~\ref{lem.lgv.signReversing.sign})
while preserving the weight $w(\mathbf{p})$
(Lemma~\ref{lem.lgv.signReversing.weight}), and has no fixed points
among intersecting path tuples
(Lemma~\ref{lem.lgv.signReversing.noFixedPoints}).  Therefore the
contributions of intersecting path tuples cancel in pairs, leaving only
the non-intersecting ones.
\end{proof}

\subsection{Helper lemmas for the sign-reversing involution}

\begin{lemma}[Vertices are distinct in acyclic digraphs]
\label{lem.lgv.path-vertices-nodup}
\lean{LGV.SimpleDigraph.Path.vertices_nodup_of_acyclic}
\leanhelper
\leanok
In an acyclic digraph $D$, all vertices on a path are distinct.
\end{lemma}

\begin{proof}
\leanok
If a vertex appeared twice at positions $i < j$, we could extract the
subpath from $i$ to $j$, which forms a cycle.  By acyclicity,
this subpath has length $1$, so $i = j$.
\end{proof}

\begin{lemma}[Exchange of tails preserves weight]
\label{lem.lgv.exchangeTails.weight}
\lean{LGV.exchangeTails_weight}
\leanhelper
\leanok
Let $D$ be an acyclic digraph with arc weight function $w$.
Let $p$ and $q$ be two paths sharing a vertex $v$.  Let $(p', q')$
be the result of exchanging the tails of $p$ and $q$ at $v$
(i.e., $p'$ takes the head of $p$ up to $v$ followed by the tail of $q$
from $v$, and symmetrically for $q'$).  Then
$w(p) \cdot w(q) = w(p') \cdot w(q')$.
\end{lemma}

\begin{proof}
The weight of a path is the product of the weights of its arcs.
When we exchange tails at $v$, arcs are redistributed between the two
paths, but the total multiset of arcs is preserved.  Hence the product
is unchanged.
\end{proof}

\begin{lemma}[Exchange of tails is an involution]
\label{lem.lgv.exchangeTails.involutive}
\lean{LGV.exchangeTails_involutive}
\leanhelper
\leanok
Let $D$ be an acyclic digraph, and let $p, q$ be paths sharing a
vertex $v$.  Let $(p', q')$ be the result of exchanging the tails of $p$ and $q$ at $v$.
Then exchanging the tails of $p'$ and $q'$ at $v$ returns $(p, q)$.
\end{lemma}

\begin{proof}
\leanok
Exchanging tails at $v$ twice returns each path to its original form,
since the head and tail of each path at $v$ are restored.
\end{proof}

\begin{lemma}[Sign-reversing map preserves intersection]
\label{lem.lgv.signReversing.isIntersecting}
\lean{LGV.signReversing_isIntersecting}
\leanhelper
\leanok
Let $D$ be a path-finite acyclic digraph.  If $(\sigma, \mathbf{p})$
is an intersecting path tuple with permutation, then
$\varphi(\sigma, \mathbf{p})$ is also intersecting.
\end{lemma}

\begin{proof}
\leanok
By Lemma~\ref{lem.lgv.exchangeTails.involutive}, exchanging tails
is an involution.  If the result were non-intersecting, applying
$\varphi$ again would give a non-intersecting tuple equal to the
original (contradicting the hypothesis that it is intersecting).
\end{proof}

\begin{lemma}[Canonical intersection data is invariant]
\label{lem.lgv.signReversing.canonical-eq}
Let $D$ be an acyclic digraph and $(\sigma, \mathbf{p})$ an
intersecting path tuple with canonical intersection data $(i, j, v)$.
After applying the sign-reversing map, the canonical intersection data
of the resulting tuple $\varphi(\sigma, \mathbf{p})$ is still $(i, j, v)$.
\end{lemma}

\begin{proof}
The head of path $i$ (vertices before $v$, inclusive) is unchanged by
the tail exchange, so $v$ remains the first crowded vertex on path $i$.
Furthermore, $i$ is still the smallest crowded path index (indices
below $i$ have unchanged paths), and $j$ is still the largest index
sharing vertex $v$ with $i$ (the set of path indices containing $v$
is preserved).
\end{proof}

\begin{lemma}[Sign-reversing involution is an involution]
\label{lem.lgv.signReversing.involutive}
\lean{LGV.signReversing_involutive}
\leanhelper
\leanok
Let $D$ be a path-finite acyclic digraph.  Let $(\sigma, \mathbf{p})$
be an intersecting path tuple with permutation.  Then the sign-reversing
map $\varphi$ satisfies: if $\varphi(\sigma, \mathbf{p})$ is
intersecting, then $\varphi(\varphi(\sigma, \mathbf{p})) = (\sigma, \mathbf{p})$.
\end{lemma}

\begin{proof}
\leanok
After applying $\varphi$, the canonical intersection data $(i, j, v)$ is
preserved (Lemma~\ref{lem.lgv.signReversing.canonical-eq}).
Applying $\varphi$ again thus uses
the same $(i, j, v)$, and the paths and permutation are restored by
Lemma~\ref{lem.lgv.exchangeTails.involutive} and the identity
$(\sigma \cdot (i\; j)) \cdot (i\; j) = \sigma$.
\end{proof}

\begin{lemma}[Sign-reversing involution reverses sign]
\label{lem.lgv.signReversing.sign}
\lean{LGV.signReversing_sign}
\leanhelper
\leanok
The sign-reversing map $\varphi$ satisfies: if
$\varphi(\sigma,\mathbf{p}) = (\sigma', \mathbf{p}')$, then
$(-1)^{\sigma'} = -(-1)^{\sigma}$.
\end{lemma}

\begin{proof}
\leanok
The permutation component of $\varphi(\sigma, \mathbf{p})$ is
$\sigma \cdot (i\; j)$ for distinct $i \ne j$
(Lemma~\ref{lem.lgv.signReversing.perm}).  Since transpositions
have sign $-1$, the sign of the new permutation is $-(-1)^\sigma$.
\end{proof}

\begin{lemma}[Sign-reversing involution preserves weight]
\label{lem.lgv.signReversing.weight}
\lean{LGV.signReversing_weight}
\leanhelper
\leanok
The sign-reversing map $\varphi$ satisfies: if
$\varphi(\sigma, \mathbf{p}) = (\sigma', \mathbf{p}')$, then
$w(\mathbf{p}') = w(\mathbf{p})$.
\end{lemma}

\begin{proof}
\leanok
Only paths $i$ and $j$ are modified (their tails are exchanged), and
the remaining paths are unchanged
(Lemma~\ref{lem.lgv.signReversing.paths}).  By
Lemma~\ref{lem.lgv.exchangeTails.weight}, the product of the weights of
paths $i$ and $j$ is preserved.  Hence the total weight
$w(\mathbf{p}) = \prod_l w(p_l)$ is unchanged.
\end{proof}

\begin{lemma}[Sign-reversing involution has no fixed points]
\label{lem.lgv.signReversing.noFixedPoints}
\lean{LGV.signReversing_no_fixed_points}
\leanhelper
\leanok
The sign-reversing map $\varphi$ has no fixed points among
intersecting path tuples: if $\mathbf{p}$ is intersecting, then
$\varphi(\sigma, \mathbf{p}) \ne (\sigma, \mathbf{p})$.
\end{lemma}

\begin{proof}
\leanok
The permutation component changes by a non-trivial transposition, so it
cannot equal $\sigma$.
\end{proof}

\begin{lemma}[Permutation of sign-reversing map]
\label{lem.lgv.signReversing.perm}
\lean{LGV.signReversing_perm}
\leanhelper
\leanok
The sign-reversing map $\varphi$ applied to $(\sigma, \mathbf{p})$
produces a permutation $\sigma' = \sigma \cdot (i\; j)$ for some
distinct indices $i \ne j$.
\end{lemma}

\begin{proof}
\leanok
By definition: the canonical intersection data selects indices $i$ and $j$,
and $\varphi$ composes $\sigma$ with the transposition $(i\; j)$.
\end{proof}

\begin{lemma}[Paths of sign-reversing map]
\label{lem.lgv.signReversing.paths}
\lean{LGV.signReversing_paths}
\leanhelper
\leanok
The sign-reversing map $\varphi$ applied to $(\sigma, \mathbf{p})$
produces paths $\mathbf{p}'$ such that: paths at indices $i$ and $j$
have their tails exchanged at the crowded vertex $v$, while all other
paths remain unchanged.
\end{lemma}

\begin{proof}
\leanok
By definition: $\varphi$ selects the canonical intersection data
$(i, j, v)$ and exchanges the tails of paths
$i$ and $j$ at $v$, leaving all other paths untouched.
\end{proof}

\subsection{Infrastructure for the LGV proof}

\begin{lemma}[Signed sum over intersecting path tuples is zero]
\label{lem.lgv.sum-ipat-zero}
\lean{LGV.sum_ipatWithPerm_signed_weight_eq_zero}
\leanhelper
\leanok
For any path-finite acyclic digraph $D$ with arc weights $w$,
\[
  \sum_{\substack{(\sigma, \mathbf{p}) \\ \mathbf{p} \text{ intersecting}}}
  (-1)^\sigma \, w(\mathbf{p}) = 0.
\]
\end{lemma}

\begin{proof}
Apply the sign-reversing involution $\varphi$ as a fixed-point-free involution.
Each pair $\{(\sigma, \mathbf{p}),\; \varphi(\sigma, \mathbf{p})\}$
contributes $(-1)^\sigma w(\mathbf{p}) + (-(-1)^\sigma) w(\mathbf{p}) = 0$.
\end{proof}

\begin{lemma}[Determinant equals sum over all path tuples with permutation]
\label{lem.lgv.det-eq-sum-all}
\lean{LGV.det_eq_sum_allPathTupleWithPerm}
\leanhelper
\leanok
For any path-finite digraph $D$ with arc weights $w$,
\[
  \det(M) = \sum_{(\sigma, \mathbf{p})} (-1)^\sigma \, w(\mathbf{p}),
\]
where the sum ranges over all pairs of a permutation $\sigma$ and a
path tuple $\mathbf{p}$ from $\mathbf{A}$ to $\sigma(\mathbf{B})$.
\end{lemma}

\begin{proof}
Expand $\det(M)$ by the Leibniz formula.  Each summand
$(-1)^\sigma \prod_i M_{i,\sigma(i)}$ expands (using the definition
$M_{i,j} = \sum_{p : A_i \to B_j} w(p)$ and distributivity) into
a sum over all path tuples from $\mathbf{A}$ to $\sigma(\mathbf{B})$.
Combine over all $\sigma$ to obtain the sum over
all pairs $(\sigma, \mathbf{p})$.
\end{proof}

\begin{lemma}[Nipat sum over permutations equals weighted nipat sum]
\label{lem.lgv.sum-nipat-eq}
\lean{LGV.sum_nipatWithPerm_eq_sum_nipatWeightSum}
\leanhelper
\leanok
The sum over non-intersecting path tuples with permutation equals
\[
  \sum_{\sigma \in S_k} (-1)^\sigma
  \sum_{\substack{\mathbf{p} \text{ is a nipat} \\ \text{from } \mathbf{A} \text{ to } \sigma(\mathbf{B})}}
  w(\mathbf{p}).
\]
\end{lemma}

\begin{proof}
\leanok
Decompose the set of non-intersecting path tuples with permutation as a disjoint union over
permutations $\sigma$ and nipats from $\mathbf{A}$ to $\sigma(\mathbf{B})$,
then factor out $(-1)^\sigma$ from each inner sum.
\end{proof}

\begin{lemma}[Product of path weight sums equals sum over path tuples]
\label{lem.lgv.prod-pathWeightSum}
\lean{LGV.prod_pathWeightSum_eq_sum_pathTupleWeight}
\leanhelper
\leanok
For a path-finite digraph $D$ with arc weights $w$,
\[
  \prod_{j=1}^{k} \biggl(\sum_{p : A_j \to B_j} w(p)\biggr)
  = \sum_{\substack{\mathbf{p} \text{ is a path tuple} \\ \text{from } \mathbf{A} \text{ to } \mathbf{B}}}
    w(\mathbf{p}).
\]
\end{lemma}

\begin{proof}
Expand the finite product of finite sums using distributivity.
The resulting index set bijects with all path tuples (choosing one
path for each component).
\end{proof}

\subsection{The nonpermutable case}

\begin{lemma}[Discrete intermediate value theorem]
\label{lem.lgv.discrete-ivt}
\lean{LGV.discrete_ivt}
\leanhelper
\leanok
Let $f : \mathbb{Z} \to \mathbb{Z}$ satisfy $|f(n+1) - f(n)| \le 1$
for all $n \in [a, b)$.  If $f(a) \ge 0$ and $f(b) \le 0$, then
there exists $c \in [a, b]$ with $f(c) = 0$.
\end{lemma}

\begin{proof}
\leanok
By induction on $b - a$.  If $f(a) = 0$, take $c = a$.  Otherwise,
$f(a) > 0$ and the step bound gives $f(a+1) \ge 0$.  If $f(a+1) \le 0$,
take $c = a + 1$; otherwise apply the induction hypothesis to $[a+1, b]$.
\end{proof}

\begin{lemma}[Baby Jordan curve theorem for lattice paths]
\label{lem.lgv.baby-jordan}
\lean{LGV.baby_jordan}
\leanhelper
\leanok
Let $A, A', B, B' \in \mathbb{Z}^2$ with
$\operatorname{x}(A') \le \operatorname{x}(A)$,
$\operatorname{y}(A) \le \operatorname{y}(A')$,
$\operatorname{x}(B') \le \operatorname{x}(B)$, and
$\operatorname{y}(B) \le \operatorname{y}(B')$.
If $p$ is a path from $A$ to $B'$ and $p'$ is a path from $A'$ to $B$,
then $p$ and $p'$ share a vertex.
\end{lemma}

\begin{proof}
Consider the ``coordinate sum'' $s = x + y$ along each path.
The range of $s$-values for $p$ is $[A_x + A_y,\; B'_x + B'_y]$,
and for $p'$ it is $[A'_x + A'_y,\; B_x + B_y]$.
These ranges overlap.  At the low end of the overlap, the
$x$-coordinate of $p$ is $\ge$ that of $p'$; at the high end, it is
$\le$.  Since the $x$-coordinate changes by at most $1$ per step,
a discrete intermediate value argument
(Lemma~\ref{lem.lgv.discrete-ivt}) yields a value of $s$ where
$p$ and $p'$ have the same $x$-coordinate.  Since they also have the
same coordinate sum $s$, the $y$-coordinates agree as well, so they
share that vertex.
\end{proof}

\begin{lemma}[Monotone permutation is identity]
\label{lem.lgv.monotone-perm-eq-id}
\lean{LGV.monotone_perm_eq_id}
\leanhelper
\leanok
If $\sigma \in S_k$ satisfies $\sigma(i) \le \sigma(j)$ whenever $i < j$,
then $\sigma = \operatorname{id}$.
\end{lemma}

\begin{proof}
\leanok
By strong induction.  If $\sigma$ is monotone, then for the smallest
element $i$ we have $\sigma(i) \ge i$ (otherwise injectivity is violated)
and $\sigma(i) \le i$ (otherwise the preimage of $i$ would have to be
smaller, violating monotonicity).  Hence $\sigma(i) = i$, and the claim
follows by induction.
\end{proof}

\begin{lemma}[No nipats for non-identity permutations under sorting]
\label{lem.lgv.no-nipats-nonid}
\lean{LGV.no_nipats_nonidentity}
\leanhelper
\leanok
Under the sorting conditions of Corollary~\ref{cor.lgv.kpaths.wt-np},
if $\sigma \in S_k$ is not the identity permutation, then there are no
non-intersecting path tuples from $\mathbf{A}$ to $\sigma(\mathbf{B})$.
\end{lemma}

\begin{proof}
\leanok
Since $\sigma \ne \operatorname{id}$, by
Lemma~\ref{lem.lgv.monotone-perm-eq-id} there exists $i \in [k-1]$
with $\sigma(i) > \sigma(i+1)$.  The sorting conditions then satisfy
the hypotheses of Lemma~\ref{lem.lgv.baby-jordan} for paths
$p_i : A_i \to B_{\sigma(i+1)}$ and $p_{i+1} : A_{i+1} \to B_{\sigma(i)}$,
so these two paths must intersect.  Hence no nipat exists.
\end{proof}

\begin{corollary}[LGV lemma, nonpermutable lattice weight version]
\label{cor.lgv.kpaths.wt-np}
\lean{LGV.lgv_nonpermutable}
\leantarget
\leanok
Consider the setting of Theorem~\ref{thm.lgv.kpaths.wt}, but
additionally assume that
\begin{align}
  \operatorname{x}(A_1) &\ge \operatorname{x}(A_2) \ge \cdots \ge \operatorname{x}(A_k);
    \label{eq.cor.lgv.kpaths.wt-np.xA} \\
  \operatorname{y}(A_1) &\le \operatorname{y}(A_2) \le \cdots \le \operatorname{y}(A_k);
    \label{eq.cor.lgv.kpaths.wt-np.yA} \\
  \operatorname{x}(B_1) &\ge \operatorname{x}(B_2) \ge \cdots \ge \operatorname{x}(B_k);
    \label{eq.cor.lgv.kpaths.wt-np.xB} \\
  \operatorname{y}(B_1) &\le \operatorname{y}(B_2) \le \cdots \le \operatorname{y}(B_k).
    \label{eq.cor.lgv.kpaths.wt-np.yB}
\end{align}
Here, $\operatorname{x}(P)$ and $\operatorname{y}(P)$ denote the two
coordinates of any point $P \in \mathbb{Z}^2$.

Then, there are no nipats from $\mathbf{A}$ to $\sigma(\mathbf{B})$
when $\sigma \in S_k$ is not the identity permutation
$\operatorname{id} \in S_k$.  Therefore, the claim of
Theorem~\ref{thm.lgv.kpaths.wt} simplifies to
\begin{equation}
  \det\!\left(\left(\sum_{p : A_i \to B_j} w(p)\right)_{1 \le i \le k,\; 1 \le j \le k}\right)
  = \sum_{\substack{\mathbf{p} \text{ is a nipat} \\ \text{from } \mathbf{A} \text{ to } \mathbf{B}}}
    w(\mathbf{p}).
  \label{eq.cor.lgv.kpaths.wt-np.claim}
\end{equation}
\end{corollary}

\begin{proof}
Let $\sigma \in S_k$ be a permutation that is not the identity
permutation $\operatorname{id} \in S_k$.
By Lemma~\ref{lem.lgv.no-nipats-nonid}, there are no nipats from
$\mathbf{A}$ to $\sigma(\mathbf{B})$, so
\[
  \sum_{\substack{\mathbf{p} \text{ is a nipat} \\ \text{from } \mathbf{A} \text{ to } \sigma(\mathbf{B})}}
    w(\mathbf{p}) = 0.
\]
Now, Theorem~\ref{thm.lgv.kpaths.wt} yields
\begin{align*}
  &\det\!\left(\left(\sum_{p : A_i \to B_j} w(p)\right)_{1 \le i \le k,\; 1 \le j \le k}\right) \\
  &= \sum_{\sigma \in S_k} (-1)^{\sigma}
    \sum_{\substack{\mathbf{p} \text{ is a nipat} \\ \text{from } \mathbf{A} \text{ to } \sigma(\mathbf{B})}}
    w(\mathbf{p}) \\
  &= \underbrace{(-1)^{\operatorname{id}}}_{=1}
    \sum_{\substack{\mathbf{p} \text{ is a nipat} \\ \text{from } \mathbf{A} \text{ to } \mathbf{B}}}
    w(\mathbf{p})
    + \sum_{\substack{\sigma \in S_k \\ \sigma \ne \operatorname{id}}} (-1)^{\sigma}
      \underbrace{\sum_{\substack{\mathbf{p} \text{ is a nipat} \\ \text{from } \mathbf{A} \text{ to } \sigma(\mathbf{B})}}
        w(\mathbf{p})}_{= 0} \\
  &= \sum_{\substack{\mathbf{p} \text{ is a nipat} \\ \text{from } \mathbf{A} \text{ to } \mathbf{B}}}
    w(\mathbf{p}).
\end{align*}
\end{proof}

\subsection{Applications}

\subsubsection{Nonnegativity of binomial coefficient determinants}

\begin{lemma}[FKG inequality for binomial coefficients]
\label{lem.lgv.binom-fkg}
\lean{LGV.binom_2x2_ineq}
\leanhelper
\leanok
Let $a, b, c, d$ be nonnegative integers with $a \ge b$ and $c \ge d$.  Then
\[
  \binom{a}{d} \binom{b}{c} \le \binom{a}{c} \binom{b}{d}.
\]
\end{lemma}

\begin{proof}
\leanok
By induction on $c - d$.  The base case $c = d$ is trivial.
For the inductive step, use the identity
$\binom{n}{m+1}(m+1) = \binom{n}{m}(n - m)$ to reduce the $(c+1)$
case to the $c$ case, together with the inequality $b - m \le a - m$
(from $b \le a$).
\end{proof}

\begin{lemma}[Lattice path count equals binomial coefficient]
\label{lem.lgv.paths-count-eq-choose}
\lean{LGV.paths_count_eq_choose}
\leanhelper
\leanok
The number of lattice paths in $\mathbb{Z}^2$ from $(0, -a)$ to
$(b, -b)$ equals $\binom{a}{b}$.
\end{lemma}

\begin{proof}
\leanok
A path from $(0,-a)$ to $(b,-b)$ has exactly $a$ steps (the coordinate
sum $x+y$ increases from $-a$ to $0$), of which exactly $b$ are
east-steps (the $x$-coordinate increases from $0$ to $b$).  Choosing
which $b$ of the $a$ steps are east-steps gives a bijection with
$b$-element subsets of $[a]$, so the count is $\binom{a}{b}$.
\end{proof}

\begin{lemma}[Lattice paths biject to subsets]
\label{lem.lgv.paths-equiv-powersetCard}
\lean{LGV.paths_count_eq_choose, LGV.paths_equiv_powersetCard}
\leanhelper
\leanok
There is a bijection between paths from $(0, -a)$ to $(b, -b)$ in the
integer lattice and $b$-element subsets of $\{1, \ldots, a\}$.
\end{lemma}

\begin{proof}
\leanok
A path is uniquely determined by the positions of its east-steps among
all $a$ steps.  This gives a bijection to $b$-element subsets of
$\{1, \ldots, a\}$.
\end{proof}

\subsubsection{Helpers for Corollary~\ref{cor.lgv.binom-det-nonneg}}

\begin{lemma}[Binomial matrix equals path weight matrix]
\label{lem.lgv.binom-matrix-eq}
\lean{LGV.binom_matrix_eq_pathWeightMatrix}
\leanhelper
\leanok
With $A_i = (0, -a_i)$, $B_j = (b_j, -b_j)$, and unit arc weights, the
path weight matrix equals the matrix $\bigl(\binom{a_i}{b_j}\bigr)$.
\end{lemma}

\begin{proof}
\leanok
The $(i,j)$ entry of the path weight matrix (with unit weights) is
the number of paths from $A_i$ to $B_j$, which by
Lemma~\ref{lem.lgv.paths-count-eq-choose} equals $\binom{a_i}{b_j}$.
\end{proof}

\begin{lemma}[Signed ipat count is zero (unit weight case)]
\label{lem.lgv.sum-signed-ipat-card-zero}
\lean{LGV.sum_signed_ipatFinset_card_eq_zero}
\leanhelper
\leanok
With unit arc weights on the integer lattice,
\[
  \sum_{\sigma \in S_k} (-1)^\sigma \,
  \bigl|\{\mathbf{p} \text{ ipat from } \mathbf{A} \text{ to } \sigma(\mathbf{B})\}\bigr| = 0.
\]
\end{lemma}

\begin{proof}
This is the specialization of the signed ipat cancellation
(Lemma~\ref{lem.lgv.sum-ipat-zero}) to unit weights, where the weight of
each path tuple is $1$ and the sum reduces to a signed cardinality count.
\end{proof}

\begin{corollary}
\label{cor.lgv.binom-det-nonneg}
\lean{LGV.binom_det_nonneg}
\leantarget
\leanok
Let $k \in \mathbb{N}$.  Let $a_1, a_2, \ldots, a_k$ and
$b_1, b_2, \ldots, b_k$ be nonnegative integers such that
\[
  a_1 \ge a_2 \ge \cdots \ge a_k
  \qquad \text{and} \qquad
  b_1 \ge b_2 \ge \cdots \ge b_k.
\]
Then,
\[
  \det\!\left(\left(\binom{a_i}{b_j}\right)_{1 \le i \le k,\; 1 \le j \le k}\right) \ge 0.
\]
\end{corollary}

\begin{proof}
Set $K = \mathbb{Z}$ and $w(a) := 1$ for each arc $a$ of $\mathbb{Z}^2$.
Define the lattice points
\[
  A_i := (0, -a_i) \qquad \text{and} \qquad B_i := (b_i, -b_i)
\]
for all $i \in [k]$.  These lattice points satisfy the assumptions
of Corollary~\ref{cor.lgv.kpaths.wt-np}.  Hence,
\eqref{eq.cor.lgv.kpaths.wt-np.claim} entails
\[
  \det\!\left(\left(\text{\# of paths from } A_i \text{ to } B_j\right)_{1 \le i \le k,\; 1 \le j \le k}\right)
  = \text{\# of nipats from } \mathbf{A} \text{ to } \mathbf{B}.
\]
Using Proposition~\ref{prop.lgv.1-paths.ct}, the matrix on the left
hand side is
$\bigl(\binom{a_i}{b_j}\bigr)_{1 \le i \le k,\; 1 \le j \le k}$.
The right hand side is $\ge 0$, so the determinant is $\ge 0$.

In Lean, the cases $k = 0$ and $k = 1$ are proved directly (empty and
$1\times 1$ determinants), $k = 2$ uses the FKG inequality
(Lemma~\ref{lem.lgv.binom-fkg}), and $k \ge 3$ uses the full LGV
machinery.
\end{proof}

\subsubsection{Catalan Hankel determinant}

\begin{lemma}[Dyck digraph is acyclic]
\label{lem.lgv.dyck-acyclic}
\lean{LGV.dyckDigraph_acyclic}
\leanhelper
\leanok
The Dyck digraph (with vertex set $\mathbb{Z} \times \mathbb{N}$ and arcs
$(i,j) \to (i+1, j+1)$ and $(i,j) \to (i+1, j-1)$ for $j > 0$) is acyclic.
\end{lemma}

\begin{proof}
\leanok
Every arc increases the first coordinate by $1$, so any path is strictly
increasing in the first coordinate.  Hence there can be no cycle.
\end{proof}

\begin{lemma}[Dyck digraph is path-finite]
\label{lem.lgv.dyck-pathfinite}
\lean{LGV.dyckDigraph_pathFinite}
\leanhelper
\leanok
The Dyck digraph is path-finite.
\end{lemma}

\begin{proof}
\leanok
Since the first coordinate increases by exactly $1$ on each arc, a path of
length $n$ visits vertices whose first coordinates form a consecutive
sequence.  The second coordinate is bounded by the first coordinate
difference, so there are only finitely many possible vertex sequences.
\end{proof}

\begin{lemma}[Dyck paths count equals Catalan number]
\label{lem.lgv.dyck-paths-eq-catalan}
\lean{LGV.dyck_paths_eq_catalan}
\leanhelper
\leanok
The number of paths from $(0, 0)$ to $(2n, 0)$ in the Dyck digraph
equals the Catalan number $c_n$.
\end{lemma}

\begin{proof}
\leanok
A path from $(0,0)$ to $(2n, 0)$ in the Dyck digraph consists of $2n$
steps, each going either up-right or down-right, and never going below
the $x$-axis.  Such a path is exactly a Dyck path of semilength~$n$.
The bijection with Dyck words (Lemma~\ref{lem.lgv.dyck-path-bijection})
establishes the correspondence with the Catalan number.
\end{proof}

\subsubsection{Helpers for the Dyck path--DyckWord bijection}

\begin{lemma}[Dyck path to DyckWord bijection]
\label{lem.lgv.dyck-path-bijection}
\lean{LGV.dyckPathToWord, LGV.dyckWordToPath}
\leanhelper
\leanok
There is a bijection between paths from $(0, 0)$ to $(2n, 0)$ in
the Dyck digraph and Dyck words of semilength $n$.
The forward direction converts each arc into an up or down step;
the backward direction constructs a vertex list from a Dyck word.
\end{lemma}

\begin{proof}
\leanok
The forward map sends each arc $(x, y) \to (x+1, y+1)$ to an up-step
and each arc $(x, y) \to (x+1, y-1)$ to a down-step, verifying that
the resulting list of steps forms a well-formed Dyck word.  The backward
map constructs vertices by scanning the Dyck word and accumulating
$y$-coordinates.  Round-trip equalities (the two maps are mutual inverses)
are verified by induction on the step list.
\end{proof}

\begin{lemma}[Shifted Dyck paths count equals Catalan number]
\label{lem.lgv.dyck-paths-shifted}
\lean{LGV.dyck_paths_shifted_eq_catalan}
\leanhelper
\leanok
The number of paths from $(-2i, 0)$ to $(2j, 0)$ in the Dyck digraph
equals $c_{i+j}$.
\end{lemma}

\begin{proof}
\leanok
By translation invariance
(Lemma~\ref{lem.lgv.translatePath-bijection}), the set of paths
from $(-2i, 0)$ to $(2j, 0)$ bijects with paths from $(0, 0)$
to $(2(i+j), 0)$, which has cardinality $c_{i+j}$ by
Lemma~\ref{lem.lgv.dyck-paths-eq-catalan}.
\end{proof}

\begin{lemma}[Translation is a bijection on Dyck paths]
\label{lem.lgv.translatePath-bijection}
For any integer $d$, translation by $(d, 0)$ is a bijection on paths
in the Dyck digraph: if $p$ is a path from $A$ to $B$, then translating
each vertex by $(d, 0)$ gives a path from $A + (d, 0)$ to $B + (d, 0)$.
The map is injective and surjective (the inverse translates by $(-d, 0)$).
\end{lemma}

\begin{proof}
Each arc $(x, y) \to (x+1, y \pm 1)$ is preserved by horizontal translation,
so the translated vertex list forms a valid path.  Injectivity follows
from injectivity of the translation map on vertices.
Surjectivity follows from the inverse translation.
\end{proof}

\begin{lemma}[Catalan Hankel matrix equals path weight matrix]
\label{lem.lgv.catalan-hankel-eq-pwm}
\lean{LGV.catalanHankelMatrix_eq_pathWeightMatrix}
\leanhelper
\leanok
The Catalan Hankel matrix $(c_{i+j})_{0 \le i,j < k}$ equals the
path weight matrix of the Dyck digraph with unit arc weights and
$k$-vertices $A_i = (-2i, 0)$, $B_j = (2j, 0)$.
\end{lemma}

\begin{proof}
\leanok
The $(i,j)$ entry of the path weight matrix (with unit weights) counts
paths from $(-2i, 0)$ to $(2j, 0)$, which equals $c_{i+j}$ by
Lemma~\ref{lem.lgv.dyck-paths-shifted}.
\end{proof}

\subsubsection{Helpers for the unique nipat (Catalan case)}

\begin{lemma}[Nested Dyck path construction]
\label{lem.lgv.nestedDyckPath}
\lean{LGV.nestedDyckPath, LGV.nestedDyckPath_start, LGV.nestedDyckPath_finish}
\leanhelper
\leanok
For each $i \in \mathbb{N}$, define the \emph{nested Dyck path}
$\gamma_i$ as the path in the Dyck digraph from $(-2i, 0)$ to $(2i, 0)$
that goes straight up (northeast) from $(-2i, 0)$ to $(0, 2i)$ and
then straight down (southeast) from $(0, 2i)$ to $(2i, 0)$.
This path has $4i$ arcs and visits the vertices
$(j - 2i,\; \min(j,\; 4i - j))$ for $j = 0, 1, \ldots, 4i$.
\end{lemma}

\begin{proof}
The vertex list is constructed explicitly.
The start and finish are verified by evaluating the vertex list at positions
$0$ and $4i$.  Each consecutive pair of vertices forms a valid Dyck arc
(either up-right or down-right, with $y \ge 0$).
\end{proof}

\begin{lemma}[Nested Dyck paths are disjoint]
\label{lem.lgv.nestedDyckPath-disjoint}
\lean{LGV.nestedDyckPath_disjoint}
\leanhelper
\leanok
For $i \ne j$, the nested Dyck paths $\gamma_i$ and $\gamma_j$
share no vertices.
\end{lemma}

\begin{proof}
\leanok
A vertex on $\gamma_i$ at position $m$ has
$y = \min(m, 4i - m)$.  At any shared $x$-coordinate,
the corresponding positions yield different $y$-coordinates
when $i \ne j$, because the maximum height $2i$ differs.
\end{proof}

\begin{lemma}[Dyck path parity constraint]
\label{lem.lgv.dyck-path-parity}
For any path $p$ in the Dyck digraph, the $y$-coordinate at step $m$
has the same parity as $y_{\mathrm{start}} + m$, where $y_{\mathrm{start}}$ is the
$y$-coordinate of the starting vertex.  (Each arc changes $y$
by $\pm 1$, so the parity alternates.)
\end{lemma}

\begin{proof}
By induction on $m$.  Each arc changes $y$ by $\pm 1$, flipping
the parity.
\end{proof}

\begin{lemma}[Dyck path height upper bound]
\label{lem.lgv.dyck-path-y-bound}
For a path $p$ in the Dyck digraph from $(-2i, 0)$ to $(2i, 0)$,
the $y$-coordinate at step $m$ satisfies
$y_m \le \min(m,\; 4i - m)$.
\end{lemma}

\begin{proof}
Each step changes $y$ by at most $+1$, so $y_m \le m$.
Symmetrically, $y$ must return to $0$ in $4i - m$ remaining steps,
so $y_m \le 4i - m$.
\end{proof}

\begin{lemma}[Nested Dyck paths form the unique nipat]
\label{lem.lgv.catalan-unique-nipat}
\lean{LGV.catalan_unique_nipat}
\leanhelper
\leanok
Define $k$-vertices
$A_i := (-2(i-1), 0)$ and $B_i := (2(i-1), 0)$ for $i \in [k]$.
There is exactly one non-intersecting path tuple from $\mathbf{A}$ to
$\mathbf{B}$ in the Dyck digraph: the \emph{nested} Dyck paths, where
path $i$ goes from $(-2i, 0)$ up to $(0, 2i)$ and back down to $(2i, 0)$.
\end{lemma}

\begin{proof}
\textbf{Existence and non-intersection:}
The nested Dyck path for index $i$ visits only vertices $(x, y)$
with $|x| \le 2i$ and $y \le 2i$, and at the midpoint $x = 0$ it
reaches height exactly $2i$.  Two nested paths for indices $i \ne j$
are disjoint because they reach different maximum heights
(Lemma~\ref{lem.lgv.nestedDyckPath-disjoint}).

\textbf{Uniqueness:}
Any nipat must equal the nested one.  At the midpoint of path $i$
(step $2i$), the $y$-coordinate must be exactly $2i$: it is at most
$2i$ (by Lemma~\ref{lem.lgv.dyck-path-y-bound}), has the correct
parity (by Lemma~\ref{lem.lgv.dyck-path-parity}), and cannot equal
$2m$ for any $m < i$ because
that would force a shared vertex with the (inductively determined)
path $m$.  Once the midpoint heights are fixed, the entire path is
forced (it must go straight up, then straight down) by parity and the
non-intersection constraint.

For $k \ge 2$, the full proof uses strong induction on the path index.
\end{proof}

\begin{corollary}
\label{cor.lgv.catalan-hankel-det-0}
\lean{LGV.catalan_hankel_det}
\leantarget
\leanok
Let $k \in \mathbb{N}$.  Recall the Catalan numbers
$c_n = \frac{1}{n+1}\binom{2n}{n}$ for all $n \in \mathbb{N}$.  Then,
\[
  \det\!\left(\left(c_{i+j-2}\right)_{1 \le i \le k,\; 1 \le j \le k}\right)
  = \det\!\begin{pmatrix}
    c_0    & c_1    & \cdots & c_{k-1} \\
    c_1    & c_2    & \cdots & c_k     \\
    \vdots & \vdots & \ddots & \vdots  \\
    c_{k-1}& c_k    & \cdots & c_{2k-2}
  \end{pmatrix}
  = 1.
\]
\end{corollary}

\begin{proof}
We use not the lattice $\mathbb{Z}^2$, but a different digraph: the
simple digraph with vertex set $\mathbb{Z} \times \mathbb{N}$ and arcs
$(i,j) \to (i+1, j+1)$ for all $(i,j) \in \mathbb{Z} \times \mathbb{N}$
and $(i,j) \to (i+1, j-1)$ for all $(i,j) \in \mathbb{Z} \times \mathbb{P}$.
It is easy to see that this digraph is acyclic and path-finite.

Set $K = \mathbb{Z}$ and $w(a) = 1$ for each arc $a$.  Define
$A_i := (-2(i-1), 0)$ and $B_i := (2(i-1), 0)$ for $i \in [k]$.

The Catalan number $c_n$ counts the paths from $(0,0)$ to $(2n,0)$ on
this digraph (Dyck paths).  Hence $c_n$ also counts the paths from
$(i,0)$ to $(2n+i, 0)$ for any $i \in \mathbb{N}$.  By
Lemma~\ref{lem.lgv.dyck-paths-eq-catalan}, the number of paths from
$A_i$ to $B_j$ is $c_{i+j-2}$.  Hence the path weight matrix is the
Catalan Hankel matrix
(Lemma~\ref{lem.lgv.catalan-hankel-eq-pwm}).

It is not hard to show (Lemma~\ref{lem.lgv.catalan-unique-nipat}) that
there is only one nipat from $\mathbf{A}$ to $\mathbf{B}$.  Moreover,
there are no nipats from $\mathbf{A}$ to $\sigma(\mathbf{B})$ for
$\sigma \ne \operatorname{id}$ (by an argument analogous to
Corollary~\ref{cor.lgv.kpaths.wt-np}).  Hence, by
Theorem~\ref{thm.lgv.kpaths.wt-dg}, the determinant equals $1$.
\end{proof}

\subsection{Bridge between path representations}

\begin{lemma}[Lattice path conversion is a bijection]
\label{lem.lgv.latticePath-equiv}
\lean{LGV.latticePath'Equiv}
\leanhelper
\leanok
There is a bijection between step-based lattice paths
(lists of east/north steps starting at a given
point) and vertex-based digraph paths
in $\mathbb{Z}^2$ starting at that point.
The forward direction computes the list of visited vertices; the
backward direction recovers the step sequence from consecutive vertex pairs.
\end{lemma}

\begin{proof}
The forward map is injective
(Lemma~\ref{lem.lgv.toPath-injective}) and surjective
(Lemma~\ref{lem.lgv.toPath-surjective}), establishing the equivalence.
\end{proof}

\begin{lemma}[Step-to-vertex map is injective]
\label{lem.lgv.toPath-injective}
\lean{LGV.LatticePath'.toPath_injective}
\leanhelper
\leanok
Different step sequences starting from the same point produce
different vertex-based paths.
\end{lemma}

\begin{proof}
\leanok
By induction on the path length.  If two step sequences produce the
same vertex list, the second vertices must agree.  Since each step
(east or north) produces a distinct target from the same source,
the first steps must be equal.  Apply the induction hypothesis to the tails.
\end{proof}

\begin{lemma}[Step-to-vertex map is surjective]
\label{lem.lgv.toPath-surjective}
\lean{LGV.LatticePath'.toPath_surjective}
\leanhelper
\leanok
Every vertex-based path in the integer lattice starting at a
given point is the image of some step sequence.
\end{lemma}

\begin{proof}
\leanok
Convert the path's vertex list to a step list by recovering each step from consecutive
vertex pairs: if the $x$-coordinate increases by $1$, the step is east; otherwise it is north.
The round-trip identity follows by construction.
\end{proof}

\begin{lemma}[Path tuple bijection preserves intersection]
\label{lem.lgv.latticePath-tuple-intersecting}
\lean{LGV.latticePath'Tuple_isIntersecting_iff}
\leanhelper
\leanok
A step-based lattice path tuple is intersecting if and only if
the corresponding vertex-based path tuple is intersecting.
\end{lemma}

\begin{proof}
\leanok
Both representations check whether any two paths share a vertex.
Since the step-to-vertex map preserves the vertex list, a vertex $v$
appears in two step-based paths if and only if it appears in the
corresponding vertex-based paths.
\end{proof}
